\chapter*{Introduction: On the Teaching of Trigonometry}
\addcontentsline{toc}{chapter}{Introduction: On the Teaching of Trigonometry}

Trigonometry occupies a distinctive position within the mathematical curriculum. It
is often encountered as a collection of methods for measuring angles, solving
triangles, and analyzing periodic phenomena. Yet the subject also serves as an
important point of contact between several major branches of mathematics.
Geometric reasoning, algebraic symbolism, transformation theory, complex numbers,
and harmonic analysis all intersect within its domain. As a result, the manner in
which trigonometry is presented has consequences extending well beyond the
immediate content of the course. A student's first encounter with trigonometric
ideas frequently shapes their understanding of what mathematics itself is and how
mathematical knowledge is organized.

The approach adopted in this text proceeds from the view that mathematical
understanding is strengthened when individual results are situated within a
broader network of relationships. Rather than presenting formulas and techniques
as isolated objects of study, the exposition seeks to trace their development from
a comparatively small collection of underlying principles. The goal is not merely
computational proficiency, although such proficiency remains essential. It is also
to cultivate the ability to reconstruct important results from more fundamental
ideas when memory alone proves insufficient.

From this perspective, trigonometry may be understood as an extended
investigation into the interaction between transformation and invariance.
Similarity theory studies those relationships preserved under changes of scale.
The unit circle isolates directional information by fixing magnitude.
Trigonometric identities emerge from the algebraic structure of rotational
composition. Complex numbers provide a unified language for scaling and rotation,
while Fourier analysis reveals how periodic structure may be decomposed into
elementary rotational modes. The sequence of topics traditionally encountered in
trigonometry therefore possesses a deeper coherence than is sometimes apparent
from a purely procedural presentation.

Throughout the text, occasional brief excursions consider systems outside
traditional mathematics: numeral systems, calendars, linguistic structure, and
other symbolic constructions. These comparative notes are not offered for
cultural interest alone. Each illustrates, in a different setting, a structural
idea also at work in the mathematics under discussion --- compositional
structure, periodicity, constrained state spaces, or transformations acting on
invariant structure. They are intentionally brief and subordinate to the
mathematical narrative, intended to remind the reader that the search for
structure is not confined to geometry or algebra.

The chapters that follow move gradually from the geometry of triangles to the
analysis of periodic phenomena. Along the way, students will encounter ratios,
circles, vectors, matrices, complex numbers, and harmonic decompositions. These
topics are often taught as separate units. The present text treats them as
successive stages in a continuous development. The hope is that readers will come
to see trigonometry not merely as a collection of techniques, but as one of the
clearest introductions to the study of mathematical structure itself.
